\documentclass[11pt,a4paper]{article}
\usepackage[utf8]{inputenc}
\usepackage[T5]{fontenc}
\usepackage[vietnamese]{babel}
\usepackage[a4paper,margin=1.35cm,headsep=4mm,footskip=7mm]{geometry}
\usepackage{tabularx,array,enumitem,fancyhdr,lastpage,pifont,amsmath,amssymb,multicol,longtable}
\setlength{\parindent}{0pt}
\setlength{\parskip}{1.5pt}
\setlength{\headheight}{15pt}
\setlength{\columnsep}{6mm}
\raggedbottom
\pagestyle{fancy}
\fancyhf{}
\lhead{\footnotesize TRUNG TÂM TOÁN TITAN THỦ ĐỨC}
\rhead{\footnotesize ĐÁP ÁN - LỚP 2 - BUỔI 1}
\cfoot{\footnotesize Trang \thepage/\pageref{LastPage}}
\setlist[itemize]{leftmargin=6mm,itemsep=0.5mm,topsep=0.5mm,parsep=0pt}
\setlist[enumerate]{leftmargin=7mm,itemsep=0.5mm,topsep=0.5mm,parsep=0pt}
\newcommand{\lessonpart}[1]{\vspace{1.5mm}\noindent\textbf{#1}\par\vspace{0.5mm}}
\newcommand{\answerline}{\noindent\dotfill\par\vspace{0.8mm}}
\newcommand{\shortblank}[1]{\makebox[#1]{\dotfill}}
\renewcommand{\arraystretch}{1.25}
\begin{document}
\begin{center}
{\Large\textbf{ĐÁP ÁN PHIẾU BÀI TẬP - LỚP 2 - BUỔI 1}}\\[0.5mm]
{\normalsize\textbf{(Dành cho giáo viên)}}
\end{center}
\section*{PHẦN C. TRẮC NGHIỆM}
\begin{tabularx}{\linewidth}{|*{10}{X|}}\hline
1-B & 2-D & 3-A & 4-C & 5-B & 6-D & 7-A & 8-C & 9-D & 10-B\\\hline
11-A & 12-C & 13-D & 14-B & 15-C & 16-A & 17-D & 18-A & 19-C & 20-B\\\hline
21-D & 22-A & 23-C & 24-B & 25-D & 26-A & 27-C & 28-B & 29-D & 30-A\\\hline
\end{tabularx}

\section*{PHẦN D. TỰ LUẬN}
\textit{Các bài mở có thể có nhiều đáp án đúng; giáo viên ưu tiên cách tìm có hệ thống và lời giải thích của học sinh.}\\
\textbf{Bài 1.} Số tiếp theo là 15. Em nhận thấy mỗi số hơn số trước 3 (câu ``thắc mắc'' là đáp án mở).\\
\textbf{Bài 2.} Hình tiếp theo là $\square$ (hai hình $\square$, $\circ$ xen kẽ). Câu nhận thấy/thắc mắc là đáp án mở.\\
\textbf{Bài 3.} 20, 24, 28. Mỗi số hơn số đứng trước 4 đơn vị.\\
\textbf{Bài 4.} Dãy lặp sau 3 hình. Đếm 3, 6, 9, 12 (4 nhóm), còn thêm 2 hình nữa nên hình thứ 14 là $\bullet$.\\
\textbf{Bài 5.} Đáp án mở; kiểm tra quy luật rõ ràng, đủ 5 phần tử.\\
\textbf{Bài 6.} Hiệu: +1, +2, +3, +4 nên hai số tiếp theo là 16 và 22 (cộng thêm 5 rồi 6).\\
\textbf{Bài 7.} Ví dụ: 1+5+9, 1+6+8, 2+4+9, 2+5+8, 2+6+7, 3+4+8, 3+5+7, 4+5+6 (có 8 cách).\\
\textbf{Bài 8.} Đáp án mở, ví dụ: 7+7; 3+4+7; 4+4+3+3; 3+3+4+4; 7+4+3. Kiểm tra tổng bằng 14.\\
\textbf{Bài 9.} 25, 28, 52, 58, 82, 85 (6 số). Lớn nhất 85, bé nhất 25.\\
\textbf{Bài 10.} Ban đầu có 6, 7, 8, 9 hoặc 10 cái (lớn hơn 5, không quá 10): 5 khả năng.\\
\textbf{Bài 11.} 14, 23, 32, 41, 50 (5 số).\\
\textbf{Bài 12.} (1,2,6), (1,3,5), (2,3,4). Có 3 bộ.\\
\textbf{Bài 13.} ABC (An - Bình - Chi), BAC, CAB. Có 3 cách.\\
\textbf{Bài 14.} Có 8 cách: 1+1+1+1+1; một lần bước 2 (4 cách); hai lần bước 2 (3 cách).\\
\textbf{Bài 15.} (a) Mỗi trọn bộ lệnh tiến thêm 2 ô: sau 5 lần ở ô 10. (b) Sau 3 lần ở ô 6, tiến thêm 3 ô: ô 9.\\
\textbf{Bài 16.} Đáp án mở: chấp nhận các hình khác nhau; kiểm tra học sinh nêu tiêu chí khác nhau và hiểu hình chỉ xoay không phải hình mới.\\
\textbf{Bài 17.} Hai hình tam giác. Có: cắt qua trung điểm hai cạnh đối diện thì được hai hình chữ nhật.\\
\textbf{Bài 18.} 7 que và 10 que. Mỗi hình vuông thêm vào cần thêm 3 que.\\
\textbf{Bài 19.} Không đồng ý. Quy luật là cộng 5, số tiếp theo là 30.\\
\textbf{Bài 20.} Có. Với 12 que (bội của 3), người đi sau luôn lấy sao cho cả hai lượt vừa đủ 3 que (đối phương lấy 1 thì mình lấy 2, lấy 2 thì mình lấy 1). Luôn để lại 9, 6, 3, 0 que nên người đi sau thắng.\\

\section*{PHẦN E. THỬ THÁCH NHÓM}
\textbf{Mission 1.} Có 2 cặp: 2+5 và 3+4 (cặp 1+6 không có vì không có thẻ 6). Biết đủ vì lấy lần lượt thẻ 1, 2, 3 và tìm thẻ còn lại.\\
\textbf{Mission 2.} Số tiếp theo là 14 (cộng 3). Dãy mới là đáp án mở, ví dụ 1, 5, 9, 13, 17.\\
\textbf{Mission 3.} Đáp án mở, ví dụ hình vuông và hình chữ nhật: cùng 4 góc, 4 cạnh; khác nhau ở độ dài các cạnh.
\end{document}
